\section{Innledning}
I dette forsøket studeres en transistorbasert FE-forsterker (felles emitter) med fokus på hvordan kretsens oppbygning påvirker signalforsterkning og signalform. Ved hjelp av både teoretiske beregninger og praktiske målinger analyseres sammenhengen mellom inngangssignal og utgangssignal, samt hvordan arbeidspunktet påvirker forsterkerens oppførsel.
\\[1em]
Det legges særlig vekt på å undersøke spenningsforsterkning, faseforskyvning og fenomener som klipping og ikke-lineær respons. Videre studeres hvordan endringer i komponentverdier, spesielt kollektormotstanden, påvirker forsterkningen.
\\[1em]
Etter gjennomført forsøk skal vi ha god kjennskap til FE-forsterkere, samt forstå sentrale begreper som spenningsforsterkning, faseforskyvning, inngangsimpedans og utgangsimpedans.